\section{AI model integration}
\label{sec:model}

The brief's wording for full points is ``seamlessly integrated, provides real-time predictions,
and handles various input types''. The application serves nine scoring methods and one language
model from Assignment 2, takes four kinds of input (an address, a transaction of that address, a
hypothetical transaction, and a choice of method), and answers each from memory or in well under a
second.

\subsection{The methods the application serves}
\label{sec:methods}

Names, thresholds and sizes are read from the Assignment 2 repository at commit \code{aa56a26}. A
transaction is flagged by a method when it can be scored and its score reaches that method's
threshold.

{\small
\begin{longtable}{@{}L{3.4cm}L{2.0cm}L{4.2cm}L{1.3cm}L{2.7cm}@{}}
\toprule
\textbf{Method} & \textbf{Family} & \textbf{What its score measures} & \textbf{Flag at} & \textbf{File} \\
\midrule
\endhead
Robust z baseline (A1) \newline \code{baseline} & Statistics & The largest $|z|$ of F1 to F4 against the address's own median and spread. Nothing is trained & 5.0 & \fp{baseline.json} \newline 1 kB \\
Isolation Forest \newline \code{isolation_forest} & Outlier detection & How easily random trees isolate the transaction & 0.589 & \fp{isolation_forest.joblib} \newline 2.8 MB \\
Local Outlier Factor \newline \code{lof} & Outlier detection & Local density against the nearest neighbours & 1.855 & \fp{stage/C3_lof.joblib} \newline 3.3 MB \\
K-Means distance \newline \code{kmeans} & Clustering & Distance to the nearest cluster centre & 9.978 & \fp{stage/C4_kmeans.joblib} \newline 2 kB \\
DBSCAN noise \newline \code{dbscan} & Clustering & Distance to the nearest dense core & 2.821 & \fp{stage/C5_dbscan.joblib} \newline 0.3 MB \\
Regression residual \newline \code{regression} & Regression & The miss between the predicted and actual size and gap & 3.5 & \fp{regression.joblib} \newline 1.7 MB \\
Regression expected range \newline \code{regression_interval} & Regression & Distance outside the conformal expected range & 0.0 & \fp{regression_interval.joblib} \newline 0.5 MB \\
Score fusion \newline \code{fusion} & Combination & Regression and Isolation Forest together, at one flag budget & 1.0 & \fp{fusion.joblib} \newline 5.0 MB \\
TxGPT surprisal \newline \code{txgpt_method} & Language model & How surprising the transaction's tokens are to TxGPT, trained from scratch on the crawl & 12.203 & \fp{txgpt_method.joblib} \newline 17.5 MB \\
\bottomrule
\end{longtable}
}

The nine files total about 31 MB and go in the zip. TxGPT's scores for the snapshot were computed
in advance and are stored with its method, so it shows a verdict for every snapshot transaction but
not for a what-if; the interface says ``not available for a hypothetical transaction'' in its row.
The tenth model, the 152M Explorer Assistant, writes explanations and scores nothing
(Section~\ref{sec:assistant}).

\textbf{The method selector (S-01).} Every snapshot transaction already has a score and a flag
from all nine methods in \fp{scores.csv}. Choosing a method changes which column decides
\code{flagged}, \code{score}, the daily flagged counts and the chart colours. The explanation then
opens with ``Flagged by Isolation Forest.'' and lists what stands out in F1 to F4, or the two largest
differences when nothing does, exactly as Assignment 2's template does for any method. The service
makes this view with \code{dataclasses.replace(predictor, selected=method)}, which shares the loaded
tables. The selector always shows which method is active and offers ``Back to the default''.

\subsection{From input to output}
\label{sec:prepost}

The report must describe preprocessing and postprocessing. These are the stages and where each
lives.

\begin{center}
\small
\begin{tabularx}{\textwidth}{@{}L{2.6cm}Y L{4.1cm}@{}}
\toprule
\textbf{Stage} & \textbf{What happens} & \textbf{Code} \\
\midrule
Validate & Address checksum and network; transaction ID; amount, direction and time ranges & \fp{app/core/validation.py}, the schemas \\
Look up & The address's rows, from the snapshot or the live cache & \fp{services/store.py} \\
Preprocess & Amount to satoshis; the hypothetical row with the history's median shape; features F1 to F4 from strictly earlier transactions; robust z-scores against two comparison windows; z clipped to $\pm 20$; direction as two indicator columns & \fp{explorer.features}, \fp{explorer.methods.model_inputs} \\
Predict & Each saved method's \code{score}; compared with its threshold & \fp{explorer.methods.flags}, the method classes \\
Postprocess & Units people read (BTC, hours, counts); each feature's role (flag, contributing, within usual); the comparison basis; the template explanation; the limitation line on flags; NaN to null & \fp{Predictor.detail}, \fp{explorer.explain} \\
Check & No banned word in any text; the response matches its model & \code{check_outputs}, the response model \\
\bottomrule
\end{tabularx}
\end{center}

Snapshot transactions were preprocessed and scored when Assignment 2 froze its results, so a GET
only reads and postprocesses. The time-safety rule holds everywhere: a feature uses only
transactions confirmed strictly earlier at the same address, and transactions with the same time are
never each other's history.

\subsection{The what-if, single and batched}
\label{sec:whatif}

\textbf{Single (E8).} \code{Predictor.score_what_if} as it is: 0.40 s measured.

\textbf{Batched (E9).} The curve needs the same transaction at many amounts. Called one by one,
41 amounts would take about 16 s\est{}. Because same-time transactions are never each other's
history, all the hypothetical rows can share one time and go through one feature build: each sees
only the real, earlier history.

\begin{lstlisting}[style=py]
rows = [hypothetical(address, amount, direction, t, txid=f"what-if-{i:03d}")
        for i, amount in enumerate(amounts)]            # same time t, distinct ids
frame = build_features(pd.concat([earlier, pd.DataFrame(rows)]), FeatureConfig.refined())
new = frame[frame["txid"].str.startswith("what-if")].assign(period="eval")
scores = {name: method.score(new) for name, method in usable.items()}
\end{lstlisting}

Measured on the running example with seven amounts from 0.0001 to 10 BTC: 0.16 s for the batch
against 1.88 s for seven single calls, and the scores of all eight methods equal to six decimal
places. The default curve of 41 amounts, from a hundredth of the usual amount to a thousand times
it, took 0.18 s for all eight methods. On it the baseline's score falls from 5.94 at the smallest
amount to 0.45 at the usual one and rises to 9.00 at the largest: an amount far below the usual
stands out as well as one far above. The flag line of 5.0 is crossed at about 0.023 BTC on the way
up and at about 0.00001 BTC on the way down.

\begin{done}
A parity test scores 41 amounts both ways on three addresses and both directions and requires equal
scores and flags for every method. E9 answers in under 1.5 s on the PC and the Mac.
\end{done}

\subsection{Copying Assignment 2's runtime}
\label{sec:vendoring}

\code{tools/vendor_a2.py --from ../COS30049A2 --commit <sha>} does the copy, so it can be repeated
and checked.

\begin{enumerate}
  \item It checks out nothing: it reads files from the given commit with \code{git show}, so
    uncommitted work cannot leak in.
  \item It finds the modules to copy by loading the entry points in a clean interpreter
    (\code{Predictor.load()}, one \code{score_address}, one \code{score_what_if}) and listing every
    \code{explorer} and \code{lm} module that was imported. On 10 November that list was 23
    \code{explorer} modules and \code{lm.tx_tokens}: \code{amounts}, \code{assistant},
    \code{baseline}, \code{classify}, \code{clustering}, \code{esplora}, \code{explain},
    \code{external}, \code{features}, \code{fusion}, \code{inject}, \code{labelled},
    \code{methods}, \code{outliers}, \code{predict}, \code{regression}, \code{scores},
    \code{settings}, \code{snapshot}, \code{split}, \code{stage_search}, \code{stages} and
    \code{txgpt}. S-05 adds the Assistant's modules of \code{lm} the same way; S-09 adds
    \code{transactions} and the frozen-model adapter.
  \item It copies the nine method files, \fp{selection.json}, \fp{profiles.joblib}, the four
    snapshot tables and the manifest, \fp{config/settings.json}, and the two evaluation tables E11
    reads.
  \item It writes \fp{backend/A2_VENDOR.json}: the source commit, the date, and a SHA-256 for every
    copied file.
\end{enumerate}

\begin{caution}
The package keeps the name \code{explorer} and its module names. The saved models are pickles that
record each class by its module path; a renamed module would make them fail to load. No file under
\fp{backend/explorer/} or \fp{backend/lm/} is edited by hand. A test recomputes every hash against
\fp{A2_VENDOR.json}.
\end{caution}

\subsection{Modifications for the web}
\label{sec:mods}

The brief asks the report to list ``any modifications for web deployment''. Every one is in a
service around the copied package, is written up in \fp{backend/MODIFICATIONS.md}, and has a test.

{\small
\begin{longtable}{@{}L{1.2cm}L{5.4cm}L{4.4cm}L{3.0cm}@{}}
\toprule
\textbf{No.} & \textbf{Modification} & \textbf{Why} & \textbf{Test} \\
\midrule
\endhead
Mod-1 & The models and tables load once, in the application's lifespan, and are shared read-only & A load takes 0.9 s; a request must not pay it & Start-up test; health reports the counts \\
Mod-2 & A per-address cache of the scored result, filled lazily and warmed for all 28 addresses in the background after start & Repeat views answer from memory & Second request faster than the first; equal bodies \\
Mod-3 & \code{score_address}'s single result is split into timeline, a page of rows and one detail & It is up to 2 MB; a view needs a few kB & Parity: the pieces rebuild the whole \\
Mod-4 & A timeline for addresses with too little history, built from the feature table & \code{score_address} raises for them, but A1's screen lists their transactions unscored & The three short-history addresses answer 200 \\
Mod-5 & The batched what-if (Section~\ref{sec:whatif}) & 41 amounts in one feature build & Parity with single calls \\
Mod-6 & A what-if that cannot be scored answers 422 with the history counts & A1's meaning of 422; the form can say which date would work & Early-time what-if \\
Mod-7 & The deciding method can be chosen per request & S-01 & Flags equal \fp{scores.csv} for each method \\
Mod-8 & Exceptions become HTTP errors in one place: \code{Unsupported} 404, \code{InsufficientHistory} 422, \code{ValueError} 400 & One error body & One test per mapping \\
Mod-9 & The Assistant loads on first use behind a lock, writes one answer at a time off the event loop, and caches answers by a hash of the facts & An answer takes about 10 s of CPU; two at once would double it & Second request is cached; a busy request gets 429 \\
Mod-10 & Start-up checks: library versions against the pins, every model file against its hash, a named count of worker threads & A clear message, not a flood of warnings & Wrong version stops with the message \\
Mod-11 & Frozen models applied to a fetched address, with the regression's residual scales set from that address's own earlier 70\% & S-09; the route Assignment 2 used for its held-out snapshot & Parity with Assignment 2's DS10 scores on one address \\
\bottomrule
\end{longtable}
}

\subsection{The Explorer Assistant (S-05)}
\label{sec:assistant}

\begin{itemize}
  \item \textbf{What the user sees.} Every explanation starts as the template, labelled ``Template:
    fixed sentences filled in from the numbers''. A button, ``Ask the Assistant to write this'',
    appears when \code{capabilities.assistant} is true. While it writes, the panel shows a progress
    line (``Writing; this takes about 10 seconds on this computer'') and a ``Keep the template''
    button that cancels. The answer replaces the template with the label ``Written by the Explorer
    Assistant, a language model the team trained; checked by code before being shown''.
  \item \textbf{The guardrail.} Six checks from \code{explorer.assistant.check}: every number comes
    from the facts, the limitation line is present on a flag, no banned word, the comparison words
    point the right way, the verdict matches, and the answer is not empty. A failed answer is
    discarded; the template stays, with a line naming the failed checks.
  \item \textbf{Not installed.} Without the weights the button is replaced by a note: ``The
    Assistant's model file (610 MB) is not installed; the README says how to add it. Explanations
    use the template.'' Nothing else changes.
  \item \textbf{Speed.} Decoding is greedy, so the same facts always give the same answer; answers
    are cached by a SHA-256 of the facts. \code{tools/precompute_assistant.py} writes the guarded
    answers for the snapshot's flagged transactions under the baseline to a JSON file that ships in
    the zip, so a marker without the weights still reads Assistant-written explanations there. A
    what-if or a fetched address needs the weights.
  \item \textbf{Limits.} One answer at a time; a second request waits its turn through a 429 and an
    automatic retry. The request's time limit is 90 seconds, the first load included.
\end{itemize}

\subsection{The live lookup (S-09, needs CR-08)}
\label{sec:live}

\begin{enumerate}
  \item The user types a valid address that is not in the snapshot. The unsupported state explains
    the coverage as before and, when live lookup is on, adds ``Fetch this address's public history''.
  \item \code{POST /api/live/lookups} starts a background job and answers 202. The job pages
    through Esplora's confirmed history, newest first, 25 transactions a page, with Assignment 2's
    client: paced requests, back-off, and a switch to the second host when the first limits the rate.
  \item Progress events drive a dialog: pages and transactions so far against the cap, and a
    Cancel button. The fetch stops at \code{maxTx} (500 by default), so it needs at most 20 pages.
  \item The history is cleaned (\code{build_transactions}), featured (\code{build_features}),
    split at 70\% (\code{assign_periods}) and scored by the frozen methods through Assignment 2's
    adapter. TxGPT is skipped unless PyTorch is installed.
  \item The result goes into the live store; the overview opens with a ``Live data, fetched
    10 Nov 2026 14:02 UTC, the most recent 500 transactions'' banner in place of the snapshot date.
\end{enumerate}

\begin{caution}
\textbf{Limits that are part of the scope.} At most six lookups an hour and one at a time;
Blockstream allows about 700 unauthenticated requests an hour from one IP address. An address with fewer than 10
transactions or 14 days of history shows the insufficient-history state. When only the most recent
500 transactions are fetched, the oldest of them have little history behind them and are shown
unscored; the banner says so. A fetched address is public blockchain data and is treated exactly
like a snapshot address: no owner, no label, the same wording and the same limitation line. The
cache is local and E20 clears it.
\end{caution}

\subsection{What stays out}

The Elliptic and BitcoinHeist classifiers (W-04), free chat (W-05) and any retraining (W-06). The
report's model section still recaps all of Assignment 2's models, and says which ones the web
application serves and why.
